% ============================================================================
%  Results section — MADDPG routing-variant architecture study (Paper 1)
%  Requires: \usepackage{graphicx,booktabs,siunitx} (or remove siunitx columns)
%  Figures live in host_data/results/main_run/figures/  (PDF preferred).
%  Suggested: \graphicspath{{host_data/results/main_run/figures/}}
% ============================================================================
\section{Results}
\label{sec:results}

% ---------------------------------------------------------------------------
\subsection{Experimental setup}
\label{sec:setup}

We evaluate seven cooperative MADDPG routing agents on an $86$-node hierarchical
service-provider topology comprising a meshed core, a distribution tier, and
$21$ provider-edge (PE) endpoints ($7$ base stations, $7$ MEC servers, and
$7$ content servers). Only the distribution-tier switches act as learning
agents ($32$ agents); core and access nodes forward along the labelled path
selected at ingress, following an MPLS-style PE model. Each agent chooses among
$K{=}3$ pre-computed shortest paths per destination.

The seven variants factor along three architectural axes: the
\emph{critic domain} (centralised critic, \textbf{CC}, vs.\ local critic,
\textbf{LC}), the \emph{value head} (simple vs.\ duelling $Q$-network), and the
optional \emph{GNN encoder} that pre-processes all agents' observations on the
agent graph before the actor and critic. We compare against an
\textbf{EVPN-SP} baseline: the identical forwarding and queueing model driven by
deterministic shortest-path (k$=0$) tunnels, i.e.\ the learning-free classical
control plane.

Traffic is generated as access-to-access flows. We report three regimes:
(i)~\emph{uniform}, where source--destination pairs are drawn uniformly;
(ii)~\emph{hotspot}, where a fraction $w{=}0.75$ of flows are concentrated from
$\{$BS1--4, MEC1--4$\}$ onto two content servers $\{$CS1, CS2$\}$, emulating a
CDN flash crowd; and (iii)~\emph{dual-link failure}, where the two primary
inter-cluster links (S8--S13 and S8--S22, capacity $300$ each) are removed at
the start of every episode, forcing rerouted traffic to share the surviving
core capacity. Each operating point is swept over an offered-load factor in
$[0.4, 3.0]\times$ and averaged over $20$--$30$ episodes of $256$ timesteps.
Unless noted, the headline metric is the \emph{resolved packet-delivery ratio}
(PDR), defined as delivered\,/\,(delivered\,+\,dropped), which is independent of
episode length and in-transit backlog.

% ---------------------------------------------------------------------------
\subsection{Training convergence}
\label{sec:training}

All seven variants converge under the shared flow reward
(Fig.~\ref{fig:train-reward}, Fig.~\ref{fig:train-loss}); curves show the raw
per-episode signal (faint) and a $41$-episode rolling mean (bold). Two effects
are already visible at training time. First, the \emph{local-critic} variants
attain markedly higher reward than their centralised counterparts
($+6$ to $+8$ reward units) and lower packet loss. Second, the GNN encoder
delays but does not destabilise convergence (best checkpoints at epoch $240$
vs.\ $150$). Table~\ref{tab:training} summarises the final operating points and
the best-checkpoint epoch used for all subsequent evaluation.

\begin{figure}[t]
  \centering
  \includegraphics[width=0.86\linewidth]{F1_training_reward.pdf}
  \caption{Training reward (rolling mean over the raw per-episode signal). The
  local-critic family (LC) converges to a higher reward than the centralised
  family (CC); the GNN encoder mainly shifts the convergence epoch.}
  \label{fig:train-reward}
\end{figure}

\begin{figure}[t]
  \centering
  \includegraphics[width=0.86\linewidth]{F2_training_pkt_loss.pdf}
  \caption{Training packet loss (rolling mean). LC variants reach the lowest
  loss, consistent with the reward ranking.}
  \label{fig:train-loss}
\end{figure}

\begin{table}[t]
  \centering
  \caption{Training summary. All variants stopped early; the best checkpoint
  (by validation reward) is used for every evaluation in this section.}
  \label{tab:training}
  \begin{tabular}{lcccc}
    \toprule
    Variant & Final reward & Final loss (\%) & Best epoch & Stop epoch \\
    \midrule
    CC-Simple        & 20.0 & 1.21 & 150 & 270 \\
    CC-Duelling      & 20.1 & 1.08 & 150 & 270 \\
    CC-Simple-GNN    & 22.1 & 1.17 & 240 & 360 \\
    CC-Duelling-GNN  & 21.7 & 1.15 & 240 & 360 \\
    LC-Duelling      & 26.7 & 0.89 & 150 & 270 \\
    LC-Simple        & 27.4 & 0.92 & 240 & 360 \\
    \textbf{LC-Duelling-GNN} & \textbf{29.0} & \textbf{0.89} & 240 & 360 \\
    \bottomrule
  \end{tabular}
\end{table}

% ---------------------------------------------------------------------------
\subsection{Uniform load: shortest path is hard to beat}
\label{sec:uniform}

Under uniform traffic (Fig.~\ref{fig:pdr-normal}), no learned policy surpasses
the EVPN-SP baseline at any load. EVPN-SP leads with $94.8\%$ PDR at nominal
load and $88.3\%$ at $2\times$ overload; the LC variants track it within
$0.5$--$1.4$ percentage points (pp), while the CC variants trail by
$3$--$4$\,pp. The multi-metric view (Table~\ref{tab:kpi-uniform}) shows the gap
is not confined to delivery: the CC variants also incur longer paths
($\sim$8.1 vs.\ 6.5--7.0 hops), higher tail delay, and higher utilisation. On a
symmetric demand matrix the reward-maximising policy is essentially
shortest-path, so the centralised critic's tendency to deviate onto longer
alternates is pure overhead here. This explains the apparent inversion with the
training reward (Sec.~\ref{sec:training}): training is performed on uniform
traffic, where LC's near-shortest-path policy is close to optimal and therefore
scores highest --- a high uniform-load reward reflects faithful shortest-path
imitation, not a load-spreading capability.

\begin{figure}[t]
  \centering
  \includegraphics[width=0.86\linewidth]{F3_pdr_vs_load_normal.pdf}
  \caption{Resolved PDR vs.\ offered load under uniform traffic. EVPN-SP (black,
  dotted) is the upper envelope; LC variants track it closely, CC variants
  trail.}
  \label{fig:pdr-normal}
\end{figure}

\begin{figure}[t]
  \centering
  \includegraphics[width=0.86\linewidth]{F4_loss_vs_load_normal.pdf}
  \caption{True per-packet loss vs.\ load under uniform traffic (mirror of
  Fig.~\ref{fig:pdr-normal}).}
  \label{fig:loss-normal}
\end{figure}

\begin{table}[t]
  \centering
  \caption{Key performance indicators under uniform traffic at nominal
  ($1\times$) and overload ($2\times$). Best per block in \textbf{bold}.}
  \label{tab:kpi-uniform}
  \small
  \begin{tabular}{lrrrrr}
    \toprule
    Variant & PDR\,(\%) & Loss\,(\%) & $d_{95}$ & Hops & Util \\
    \midrule
    \multicolumn{6}{l}{\emph{Load $1.0\times$}}\\
    CC-Simple        & 91.4 & 8.6 & 14.1 & 8.16 & 0.016 \\
    CC-Duelling      & 91.9 & 8.1 & 14.3 & 8.14 & 0.014 \\
    LC-Duelling      & 93.7 & 6.3 & 11.2 & 6.98 & 0.013 \\
    LC-Simple        & 93.8 & 6.2 & 10.5 & 6.89 & 0.011 \\
    LC-Duelling-GNN  & 94.1 & 5.9 & 10.1 & 6.62 & 0.011 \\
    \textbf{EVPN-SP} & \textbf{94.8} & \textbf{5.2} & \textbf{9.9} & \textbf{6.53} & 0.011 \\
    \midrule
    \multicolumn{6}{l}{\emph{Load $2.0\times$}}\\
    CC-Simple        & 84.9 & 15.0 & 13.7 & 8.08 & 0.041 \\
    CC-Duelling      & 84.4 & 15.4 & 13.7 & 8.09 & 0.038 \\
    LC-Duelling      & 87.4 & 12.5 & 10.3 & 7.11 & 0.034 \\
    LC-Simple        & 87.4 & 12.5 & 10.2 & 6.98 & 0.033 \\
    LC-Duelling-GNN  & 87.8 & 12.1 & 10.1 & 6.76 & 0.031 \\
    \textbf{EVPN-SP} & \textbf{88.3} & \textbf{11.7} & \textbf{9.9} & \textbf{6.70} & 0.031 \\
    \bottomrule
  \end{tabular}
\end{table}

% ---------------------------------------------------------------------------
\subsection{Traffic concentration: where the centralised critic pays off}
\label{sec:hotspot}

The picture inverts under the hotspot matrix (Fig.~\ref{fig:pdr-hotspot}). The
CC variants now \emph{overtake} EVPN-SP, and the advantage widens with load:
CC-Simple reaches $+5.5$\,pp over the baseline at $2\times$ and
$+8.2$\,pp at $3\times$. Crucially, the LC variants do \emph{not} share this
advantage --- at high load they fall \emph{below} EVPN-SP
(LC-Duelling: $63.6\%$ vs.\ $66.5\%$ at $3\times$).
Figure~\ref{fig:gap} makes the regime dependence explicit, plotting each
variant's PDR gap against EVPN-SP for uniform (left) and hotspot (right): the CC
curves cross from negative to strongly positive, while the LC curves hug or
dip below the baseline.

The mechanism is visible in Table~\ref{tab:mechanism}. The LC variants match
EVPN-SP to within measurement noise on every structural metric --- hop count
($8.2$ vs.\ $8.1$), mean and tail utilisation ($0.045/0.112$ vs.\
$0.044/0.110$), and overload fraction --- confirming that LC has learned a
near-shortest-path policy and therefore \emph{cannot} relieve the hotspot.
The CC variants instead trade path length for spread: they route
$\sim$1 hop longer and push utilisation across more links (util-$p_{95}$
$0.140$ vs.\ $0.110$), pulling traffic off the congested egress and lifting
delivery. This is a direct consequence of the critic's information set: a
\emph{central} critic observes global link state and can credit-assign the
system-wide benefit of a longer detour, whereas a \emph{local} critic perceives
only its adjacent links and collapses to greedy shortest-path forwarding.
LC-Duelling's small uncoordinated deviations are in fact counter-productive ---
it records the \emph{highest} overload fraction ($0.721$) while still taking
longer paths than EVPN-SP, explaining why it alone dips below the baseline.

\begin{figure}[t]
  \centering
  \includegraphics[width=0.86\linewidth]{F5_pdr_vs_load_hotspot.pdf}
  \caption{Resolved PDR vs.\ load under hotspot traffic. The CC variants
  overtake EVPN-SP and the margin grows with load; LC variants fall to or below
  the baseline.}
  \label{fig:pdr-hotspot}
\end{figure}

\begin{figure*}[t]
  \centering
  \includegraphics[width=0.92\linewidth]{F6_maddpg_vs_evpn_gap.pdf}
  \caption{PDR gap of each MADDPG variant relative to EVPN-SP
  ($+$ means the learned policy is better). Left: uniform traffic, where every
  variant is at or below the baseline. Right: hotspot traffic, where the
  centralised-critic variants cross above zero and diverge with load.}
  \label{fig:gap}
\end{figure*}

\begin{table}[t]
  \centering
  \caption{Forwarding mechanism under hotspot traffic at $3\times$ load. LC
  variants are statistically indistinguishable from shortest path; CC variants
  trade longer paths for wider load spread (higher util-$p_{95}$).}
  \label{tab:mechanism}
  \small
  \begin{tabular}{lrrrr}
    \toprule
    Variant & PDR\,(\%) & Hops & Util-$p_{95}$ & Overload \\
    \midrule
    \textbf{CC-Simple}   & \textbf{74.7} & 9.18 & 0.140 & 0.649 \\
    CC-Duelling          & 73.2 & 9.47 & 0.147 & 0.676 \\
    LC-Simple            & 67.4 & 8.23 & 0.112 & 0.698 \\
    LC-Duelling          & 63.6 & 8.45 & 0.111 & 0.721 \\
    EVPN-SP              & 66.5 & 8.15 & 0.110 & 0.708 \\
    \bottomrule
  \end{tabular}
\end{table}

% ---------------------------------------------------------------------------
\subsection{Robustness to link failures}
\label{sec:failure}

We next remove the two primary inter-cluster links and repeat both traffic
regimes. Under \emph{uniform} demand (Fig.~\ref{fig:pdr-fail-uni}) the failure
ordering mirrors the no-failure case: EVPN-SP and the LC variants are the most
robust ($\sim$83\% PDR at $2\times$), while the CC variants degrade further.
The path-stretch view (Fig.~\ref{fig:hops-fail}) again explains why --- after
the failure, the CC variants reroute over noticeably longer detours
($8.4$--$8.8$ vs.\ $7.4$ hops for EVPN-SP), and on an otherwise uniform load
that extra length competes for scarce post-failure capacity.

Under the \emph{combined} failure-plus-hotspot stress
(Fig.~\ref{fig:pdr-fail-hot}) the result inverts once more, and most sharply:
CC-Simple delivers $86.9\%$ at $2\times$ against EVPN-SP's $82.4\%$
($+4.5$\,pp), the largest learned-policy advantage we observe. When the network
is both degraded \emph{and} asymmetrically loaded --- precisely the condition in
which shortest path floods the single surviving bottleneck --- the centralised
critic's load-spreading is most valuable.

\begin{figure}[t]
  \centering
  \includegraphics[width=0.86\linewidth]{F7_pdr_vs_load_failure_uniform.pdf}
  \caption{Resolved PDR under dual-link failure, uniform traffic. EVPN-SP and
  the LC variants remain the most robust.}
  \label{fig:pdr-fail-uni}
\end{figure}

\begin{figure}[t]
  \centering
  \includegraphics[width=0.86\linewidth]{F8_pdr_vs_load_failure_hotspot.pdf}
  \caption{Resolved PDR under dual-link failure, hotspot traffic. CC-Simple
  attains the largest margin over EVPN-SP of any scenario.}
  \label{fig:pdr-fail-hot}
\end{figure}

\begin{figure}[t]
  \centering
  \includegraphics[width=0.86\linewidth]{F9_hops_vs_load_failure.pdf}
  \caption{Mean delivered hop count under failure (uniform). CC variants reroute
  over longer detours than EVPN-SP and the LC variants, accounting for their
  weaker uniform-failure delivery.}
  \label{fig:hops-fail}
\end{figure}

% ---------------------------------------------------------------------------
\subsection{Effect of the GNN encoder}
\label{sec:gnn}

A recurring motivation for graph encoders in DRL routing is improved
generalisation to topology change. Our results do not support that benefit in
this setting. Figure~\ref{fig:gnn} reports the PDR change from adding a GNN to
each base variant, for the no-failure and failure conditions. The effect is
small and inconsistent in sign ($-0.3$ to $+1.3$\,pp), and in particular the GNN
does \emph{not} confer a robustness advantage under link failure. We attribute
this to an implementation property shared with much of the literature: the
message-passing adjacency is fixed at training time, so after a link failure the
encoder continues to propagate over edges that no longer exist. The
parameter-sharing benefit of the GNN is retained, but the topology-generalisation
benefit requires recomputing the graph at inference --- which a statically
compiled adjacency does not provide.

\begin{figure}[t]
  \centering
  \includegraphics[width=0.82\linewidth]{F10_gnn_paired_delta.pdf}
  \caption{PDR change from adding a GNN encoder (GNN $-$ non-GNN) at $2\times$
  load, for each base variant, with and without link failure. The effect is
  small and not consistently positive; the encoder provides no robustness gain.}
  \label{fig:gnn}
\end{figure}

% ---------------------------------------------------------------------------
\subsection{Summary}
\label{sec:summary}

Table~\ref{tab:master} consolidates delivery at $2\times$ load across all four
operating conditions. Three findings stand out. (1)~\textbf{Critic domain is the
dominant design choice}, and its optimum is regime-dependent: the local critic
matches shortest path and wins under uniform load and uniform failure, whereas
the centralised critic learns genuine load-spreading and wins whenever traffic
is concentrated, with or without failures. (2)~\textbf{No learned policy
dominates the classical baseline everywhere}; the value of learning is realised
specifically under demand asymmetry and topology degradation. (3)~\textbf{The
GNN encoder and the duelling head are second-order}, neither changing the
ranking nor delivering the expected robustness benefit. Practically, an operator
should prefer a centralised-critic agent when traffic concentration or failures
are the dominant risk, and is otherwise well served by shortest path.

\begin{table}[t]
  \centering
  \caption{Resolved PDR (\%) at $2\times$ overload across all conditions. Best
  per column in \textbf{bold}; the regime-dependent winner shifts between the LC
  family (uniform) and CC-Simple (concentrated / failed).}
  \label{tab:master}
  \small
  \begin{tabular}{lrrrr}
    \toprule
    Variant & Uniform & Hotspot & Fail-uni & Fail-hot \\
    \midrule
    CC-Simple        & 84.9 & \textbf{81.7} & 77.4 & \textbf{86.9} \\
    CC-Duelling      & 84.4 & 80.3 & 80.0 & 83.3 \\
    CC-Simple-GNN    & 84.6 & 77.3 & 78.7 & 82.9 \\
    CC-Duelling-GNN  & 84.4 & 74.2 & 79.6 & 82.9 \\
    LC-Duelling      & 87.4 & 73.1 & 82.0 & 83.5 \\
    LC-Simple        & 87.4 & 75.9 & 82.9 & 83.6 \\
    LC-Duelling-GNN  & 87.8 & 75.2 & 82.8 & 82.9 \\
    EVPN-SP          & \textbf{88.3} & 74.9 & \textbf{83.0} & 82.4 \\
    \bottomrule
  \end{tabular}
\end{table}
